% Options for packages loaded elsewhere
\PassOptionsToPackage{unicode}{hyperref}
\PassOptionsToPackage{hyphens}{url}
\documentclass[
]{article}
\usepackage{xcolor}
\usepackage{amsmath,amssymb}
\setcounter{secnumdepth}{-\maxdimen} % remove section numbering
\usepackage{iftex}
\ifPDFTeX
  \usepackage[T1]{fontenc}
  \usepackage[utf8]{inputenc}
  \usepackage{textcomp} % provide euro and other symbols
\else % if luatex or xetex
  \usepackage{unicode-math} % this also loads fontspec
  \defaultfontfeatures{Scale=MatchLowercase}
  \defaultfontfeatures[\rmfamily]{Ligatures=TeX,Scale=1}
\fi
\usepackage{lmodern}
\ifPDFTeX\else
  % xetex/luatex font selection
\fi
% Use upquote if available, for straight quotes in verbatim environments
\IfFileExists{upquote.sty}{\usepackage{upquote}}{}
\IfFileExists{microtype.sty}{% use microtype if available
  \usepackage[]{microtype}
  \UseMicrotypeSet[protrusion]{basicmath} % disable protrusion for tt fonts
}{}
\makeatletter
\@ifundefined{KOMAClassName}{% if non-KOMA class
  \IfFileExists{parskip.sty}{%
    \usepackage{parskip}
  }{% else
    \setlength{\parindent}{0pt}
    \setlength{\parskip}{6pt plus 2pt minus 1pt}}
}{% if KOMA class
  \KOMAoptions{parskip=half}}
\makeatother
\usepackage{color}
\usepackage{fancyvrb}
\newcommand{\VerbBar}{|}
\newcommand{\VERB}{\Verb[commandchars=\\\{\}]}
\DefineVerbatimEnvironment{Highlighting}{Verbatim}{commandchars=\\\{\}}
% Add ',fontsize=\small' for more characters per line
\newenvironment{Shaded}{}{}
\newcommand{\AlertTok}[1]{\textcolor[rgb]{1.00,0.00,0.00}{\textbf{#1}}}
\newcommand{\AnnotationTok}[1]{\textcolor[rgb]{0.38,0.63,0.69}{\textbf{\textit{#1}}}}
\newcommand{\AttributeTok}[1]{\textcolor[rgb]{0.49,0.56,0.16}{#1}}
\newcommand{\BaseNTok}[1]{\textcolor[rgb]{0.25,0.63,0.44}{#1}}
\newcommand{\BuiltInTok}[1]{\textcolor[rgb]{0.00,0.50,0.00}{#1}}
\newcommand{\CharTok}[1]{\textcolor[rgb]{0.25,0.44,0.63}{#1}}
\newcommand{\CommentTok}[1]{\textcolor[rgb]{0.38,0.63,0.69}{\textit{#1}}}
\newcommand{\CommentVarTok}[1]{\textcolor[rgb]{0.38,0.63,0.69}{\textbf{\textit{#1}}}}
\newcommand{\ConstantTok}[1]{\textcolor[rgb]{0.53,0.00,0.00}{#1}}
\newcommand{\ControlFlowTok}[1]{\textcolor[rgb]{0.00,0.44,0.13}{\textbf{#1}}}
\newcommand{\DataTypeTok}[1]{\textcolor[rgb]{0.56,0.13,0.00}{#1}}
\newcommand{\DecValTok}[1]{\textcolor[rgb]{0.25,0.63,0.44}{#1}}
\newcommand{\DocumentationTok}[1]{\textcolor[rgb]{0.73,0.13,0.13}{\textit{#1}}}
\newcommand{\ErrorTok}[1]{\textcolor[rgb]{1.00,0.00,0.00}{\textbf{#1}}}
\newcommand{\ExtensionTok}[1]{#1}
\newcommand{\FloatTok}[1]{\textcolor[rgb]{0.25,0.63,0.44}{#1}}
\newcommand{\FunctionTok}[1]{\textcolor[rgb]{0.02,0.16,0.49}{#1}}
\newcommand{\ImportTok}[1]{\textcolor[rgb]{0.00,0.50,0.00}{\textbf{#1}}}
\newcommand{\InformationTok}[1]{\textcolor[rgb]{0.38,0.63,0.69}{\textbf{\textit{#1}}}}
\newcommand{\KeywordTok}[1]{\textcolor[rgb]{0.00,0.44,0.13}{\textbf{#1}}}
\newcommand{\NormalTok}[1]{#1}
\newcommand{\OperatorTok}[1]{\textcolor[rgb]{0.40,0.40,0.40}{#1}}
\newcommand{\OtherTok}[1]{\textcolor[rgb]{0.00,0.44,0.13}{#1}}
\newcommand{\PreprocessorTok}[1]{\textcolor[rgb]{0.74,0.48,0.00}{#1}}
\newcommand{\RegionMarkerTok}[1]{#1}
\newcommand{\SpecialCharTok}[1]{\textcolor[rgb]{0.25,0.44,0.63}{#1}}
\newcommand{\SpecialStringTok}[1]{\textcolor[rgb]{0.73,0.40,0.53}{#1}}
\newcommand{\StringTok}[1]{\textcolor[rgb]{0.25,0.44,0.63}{#1}}
\newcommand{\VariableTok}[1]{\textcolor[rgb]{0.10,0.09,0.49}{#1}}
\newcommand{\VerbatimStringTok}[1]{\textcolor[rgb]{0.25,0.44,0.63}{#1}}
\newcommand{\WarningTok}[1]{\textcolor[rgb]{0.38,0.63,0.69}{\textbf{\textit{#1}}}}
\usepackage{longtable,booktabs,array}
\newcounter{none} % for unnumbered tables
\usepackage{calc} % for calculating minipage widths
% Correct order of tables after \paragraph or \subparagraph
\usepackage{etoolbox}
\makeatletter
\patchcmd\longtable{\par}{\if@noskipsec\mbox{}\fi\par}{}{}
\makeatother
% Allow footnotes in longtable head/foot
\IfFileExists{footnotehyper.sty}{\usepackage{footnotehyper}}{\usepackage{footnote}}
\makesavenoteenv{longtable}
\setlength{\emergencystretch}{3em} % prevent overfull lines
\providecommand{\tightlist}{%
  \setlength{\itemsep}{0pt}\setlength{\parskip}{0pt}}
\usepackage[]{biblatex}
\usepackage{bookmark}
\IfFileExists{xurl.sty}{\usepackage{xurl}}{} % add URL line breaks if available
\urlstyle{same}
\hypersetup{
  hidelinks,
  pdfcreator={LaTeX via pandoc}}

\author{}
\date{}

\begin{document}

\section{ProtBench: a unified benchmark suite for protein representation
evaluation}\label{protbench-a-unified-benchmark-suite-for-protein-representation-evaluation}

Author: Dan Ofer

Affiliation: The Hebrew University of Jerusalem

Correspondence: Dan Ofer

Repository: \url{https://github.com/ddofer/ProtBench}

\subsection{Abstract}\label{abstract}

\textbf{Motivation:} Protein language models and sequence-based protein
predictors are often evaluated with different train/test splits, task
subsets, metrics, and probe implementations. This makes it difficult to
decide whether an observed difference reflects the model or the
evaluation harness.

\textbf{Results:} ProtBench is an open benchmark suite for evaluating
protein representations under shared data handling, probe fitting, and
metric reporting. The current registry contains 57 public tasks covering
sequence-level classification and regression, multilabel function
prediction, retrieval, residue-level labelling, pairwise
residue-residue contact prediction, and ProteinGym variant-effect
prediction. It supports frozen probes, fine-tuning modes,
ProteinGym zero-shot scoring, and non-neural k-mer, MMseqs2, and phmmer
baselines. Results are written as explicit per-task CSV rows and can be
collected into a long-format table for model comparison. A shared
evaluation slice exposes divergent effects of training objectives: while
scale (ESM-2 150M) and sentence-level contrastive learning (ProtSent)
improve homology matching, contrastive training severely degrades local
fitness tasks.

\textbf{Availability and implementation:} ProtBench is implemented in
Python with Hugging Face integrations and optional specialty tools.
Source code, task definitions, and dataset provenance notes are
available at \url{https://github.com/ddofer/ProtBench}.

\subsection{1 Introduction}\label{introduction}

Protein language models such as ESM and ProtTrans models are now common
starting points for protein sequence analysis
\autocite{rives2021biological,lin2023evolutionary,elnaggar2021prottrans}.
New models are usually reported on a subset of downstream tasks, but the
evaluation details vary: papers differ in splits, metric choices,
pooling strategies, probe implementations, and whether sequence
alignment baselines are included. These differences are enough to
obscure small or task-specific gains.

ProtBench addresses this evaluation problem rather than proposing a new
model. It provides a single task registry and command-line harness for
running many protein model families through the same splits, embedding
path, probes, and metrics. The default use case is frozen-probe
evaluation: a model is used once to embed each protein, then a small
supervised probe is trained on those embeddings. This measures
information that is accessible in the pretrained representation, and it
is cheap enough to run early during model development. More expensive
fine-tuning and zero-shot variant-effect modes are available when the
scientific question requires them.

\subsection{2 Design Principles}\label{design-principles}

ProtBench enforces four principles for reproducible protein model
evaluation.

\begin{enumerate}
\def\labelenumi{\arabic{enumi}.}
\tightlist
\item
  \textbf{Evaluation contracts are explicit.} Task, split, metric,
  probe, seed, and result row are fixed and reused across models.
\item
  \textbf{Frozen probes measure accessible information.} They test what
  is already present in the representation, independent of task-specific
  adaptation. Fine-tuning and zero-shot modes answer related questions.
\item
  \textbf{Baselines clarify the contribution.} k-mer, MMseqs2, and
  phmmer separate composition, homology, and learned effects.
\item
  \textbf{Provenance is auditable.} Sources are verified, mapped, or
  flagged; results can be traced back to original benchmarks.
\end{enumerate}

\subsection{3 System and Methods}\label{system-and-methods}

ProtBench is organized around a \texttt{TaskConfig} registry in
\texttt{benchmark\_tasks.py}. Each task records the dataset source,
sequence and label columns, problem type, split handling, and primary
metric. The main benchmark runner loads a model, embeds the relevant
sequences, fits the selected probe, evaluates on a declared split, and
writes one CSV row per task, split, seed, and probe. The same registry
is reused by comparison, result-collection, fine-tuning, residue-probe,
and alignment-baseline scripts.

Linear probes measure accessible information independent of task
adaptation. k-nearest neighbour probes test whether embedding distance
preserves local structure, which is useful for retrieval and annotation
transfer. Gradient-boosted trees test non-linear separability.
Sequence-level fine-tuning supports frozen probes, full fine-tuning,
last-layer updates, and LoRA adapters \autocite{hu2021lora}.
Residue-level tasks use token embeddings and per-residue metrics.
ProteinGym zero-shot tasks use masked-marginal scoring for substitutions
and pseudo-log-likelihood differences for indels, with the clinical sign
convention chosen so higher scores correspond to the positive/pathogenic
class \autocite{notin2023proteingym}.

Three non-neural baselines are included. k-mer frequencies provide a
composition floor. MMseqs2 and phmmer provide alignment-based baselines
using the same prepared query and target splits as model probes
\autocite{steinegger2017mmseqs2,eddy2011accelerated,larralde2023pyhmmer}.
These baselines are important for tasks where sequence similarity is
already highly informative.

Contact prediction is scored two ways from primary sequence alone. A
supervised pairwise probe builds features from frozen per-residue
embeddings of residues \(i\) and \(j\) and fits a logistic model, so it
runs against any model the registry loads. A zero-shot categorical
Jacobian substitutes all twenty amino acids at each position and measures
the resulting shift in the model's output distribution elsewhere,
yielding a coupling matrix without any training; it requires a
masked-language-model head and so ships as a separate script. Both report
precision over the top \(L/k\) ranked pairs within short, medium, and
long-range sequence-separation bands, so the two are directly comparable.
Structural coordinates build the labels only and are never model input,
and no multiple-sequence alignment is used.

\subsection{4 Benchmark Composition}\label{benchmark-composition}

The registry currently spans 57 tasks: 12 binary, 5 multiclass, 5
multilabel, and 17 regression tasks, plus 1 retrieval, 16
residue-level token classification tasks, and 1 pairwise
contact-prediction task. Forty-nine are standard probe
tasks; 8 are ProteinGym evaluations. The generated task inventory in
\texttt{paper/generated/task\_inventory.tsv} records each task key,
display name, metric, preset, and dataset identifier.

Dataset provenance is tracked separately from implementation details.
Verified entries include NetSurfP-2.0 secondary structure and
missing-coordinate disorder, SignalP 6.0 signal peptide labels, and a
locally built DisProt 2024 residue-level disorder task
\autocite{klausen2019netsurfp,teufel2022signalp,aspromonte2024disprot}.
Mapped entries are linked to original benchmark papers but still need
per-task reverification; best-effort entries are third-party rehosts
requiring manual confirmation. Mapped resources include ProteinGym,
FLIP, CATH/EAT, SCOPe, DeepLoc, CAFA, TAPE/GFP fluorescence, Meltome
Atlas, PEER, ProFET, and PPI benchmark splits
\autocite{notin2023proteingym,dallago2021flip,heinzinger2022contrastive,fox2014scope,almagro2017deeploc,zhou2019cafa,rao2019tape,sarkisyan2016local,jarzab2020meltome,xu2022peer,ofer2015profet,bernett2024ppi}.
Results should cite these original sources rather than Hugging Face
mirrors. Tasks whose Hugging Face rehosts have not yet been traced to an
original paper are retained in the software but flagged in the
supplement and dataset documentation.

\subsection{5 Biological Insights from Training
Objectives}\label{biological-insights-from-training-objectives}

A unified harness reveals how different training interventions
structurally alter model representations. Table 1 compares a vanilla
ESM-2 baseline against parameter scaling (ESM-2 150M) and contrastive
sentence-transformer tuning (ProtSent-V2) \autocite{protsent2026}. All
rows use the identical frozen linear-probe test split.

\textbf{Table 1. Biological insights across model scales and training
objectives.} Higher is better. Vanilla runs originate from the
ProteinJEPA benchmark suite; ProtSent is from the ProtSent suite.

{\def\LTcaptype{none} % do not increment counter
\begin{longtable}[]{@{}llrrr@{}}
\toprule\noalign{}
Task & Metric & Vanilla 35M & Vanilla 150M & ProtSent 35M \\
\midrule\noalign{}
\endhead
\bottomrule\noalign{}
\endlastfoot
Remote homology & Accuracy & 0.647 & 0.688 & 0.702 \\
Solubility & AUC & 0.697 & 0.721 & 0.698 \\
beta-lactamase & Spearman & 0.733 & 0.811 & 0.609 \\
Variant effect & Spearman & 0.844 & 0.843 & 0.813 \\
Stability & Spearman & 0.437 & 0.704 & 0.388 \\
\end{longtable}
}

The results show performance does not move in a single direction. First,
parameter scaling (Vanilla 35M to 150M) provides massive gains on
structural regression tasks (e.g., Stability jumps from 0.437 to 0.704),
while local mutational fitness capabilities often saturate early
(Variant effect stays flat at \textasciitilde0.84)---a dynamic
previously observed in models like ProteinBERT
\autocite{brandes2022proteinbert}. Second, contrastive learning
(ProtSent) improves global structure retrieval (homology) but severely
collapses local mutational fitness (beta-lactamase, stability, variant
effect). Relying on isolated literature metrics across different
preprocessing pipelines would obscure these divergent trade-offs.

The CATH v4.3 superfamily transfer task exemplifies ProtBench's
task-design approach. Its queries have no detectable sequence-alignment
relative in the lookup set, making it a stress test for whether
embeddings preserve fold relationships beyond sequence search. ProtBench
exposes this task through the standard k-NN annotation-transfer
interface. The scorer (\texttt{cath\_levels.py}) and optional
Tucker-head training reproduction
(\texttt{train\_cath\_tucker\_head.py}) both include self-check
validation modes. New CATH accuracy claims are deferred until committed
result artifacts are available; current usage is a demonstration of task
design and reproducible scoring.

\subsection{6 Reproducibility and Use}\label{reproducibility-and-use}

ProtBench is designed to be run from a fresh clone with public data.
Most tasks download through \texttt{datasets.load\_dataset}; local tasks
such as DisProt, conservation labels, and FLIP2 subtasks are built with
scripts in \texttt{scripts/}. The command below runs one model on a
single task and writes a benchmark CSV:

\begin{Shaded}
\begin{Highlighting}[]
\ExtensionTok{python}\NormalTok{ protein\_benchmark\_suite.py }\AttributeTok{{-}m}\NormalTok{ facebook/esm2\_t6\_8M\_UR50D }\DataTypeTok{\textbackslash{}}
    \AttributeTok{{-}{-}tasks}\NormalTok{ solubility }\AttributeTok{{-}p}\NormalTok{ linear }\AttributeTok{{-}{-}eval\_split}\NormalTok{ test}
\end{Highlighting}
\end{Shaded}

For multi-model analysis, \texttt{collect\_bench\_results.py} normalizes
probe CSVs and fine-tuning JSONL outputs into
\texttt{results/bench\_results\_all.csv}, one row per model, task,
probe, split, and metric. The paper assets can be refreshed with:

\begin{Shaded}
\begin{Highlighting}[]
\ExtensionTok{python3}\NormalTok{ scripts/paper\_assets.py }\AttributeTok{{-}{-}out{-}dir}\NormalTok{ paper/generated}
\end{Highlighting}
\end{Shaded}

\subsection{7 Limitations}\label{limitations}

Frozen probes measure representation capacity before task adaptation;
they do not replace full fine-tuning. Third-party dataset rehosts may
have drifted provenance, filtering, or labels, so formal claims should
be checked against the original sources. ProteinGym zero-shot indel
scoring is slow on large assays. Task composition, homology, and
imbalance often dominate results; k-mer and alignment baselines help
isolate learned effects.

\subsection{8 Conclusion}\label{conclusion}

ProtBench provides a practical evaluation layer for protein
representation work: one task registry, shared probes and metrics,
public dataset provenance notes, alignment and composition baselines,
and reusable result tables. Its main value is not a new modelling claim,
but a lower-friction way to make protein model comparisons easier to
reproduce and easier to audit.

\subsection{References}\label{references}

See \url{references.bib}.

\printbibliography

\end{document}
